\documentclass[11pt,a4paper]{article}

\usepackage[a4paper,margin=2.2cm]{geometry}
\usepackage{booktabs}
\usepackage[table]{xcolor}
\usepackage{graphicx}
\usepackage{pdflscape}
\usepackage{longtable}
\usepackage{caption}
\usepackage{fancyhdr}
\usepackage{xurl}
\usepackage{hyperref}
\usepackage{amsmath}
\usepackage{array}
\usepackage{float}
\usepackage{placeins}

% Allow breaks inside DOI UUID tails (e.g. after "uuid:") so long
% \url links in the event-log list do not overrun the right margin.
\makeatletter
\g@addto@macro\UrlBreaks{%
  \do\-\do\:\do\/\do\.\do\_%
  \do\0\do\1\do\2\do\3\do\4\do\5\do\6\do\7\do\8\do\9%
  \do\a\do\b\do\c\do\d\do\e\do\f%
  \do\A\do\B\do\C\do\D\do\E\do\F%
}
\makeatother
\Urlmuskip=0mu plus 1mu\relax

\hypersetup{
  colorlinks=true,
  linkcolor=black,
  urlcolor=blue,
  citecolor=black,
  bookmarksdepth=3
}

\pagestyle{fancy}

\title{Online Appendix for the ICPM 2027 submission: ``Assessing Power-Law Compatibility of Process Behavior Frequencies Across Orders of Behavioral Abstraction''}
\author{}
\date{November 20, 2026}

\begin{document}
% Auto-generated by generate.py — do not edit by hand.
\newcommand{\AppendixBestAlternativeInput}{\input{"../../../results/perm/powerlaw_statistics/real/xmin-adaptive/clauset_results/min[50]-gofp[0.10]-pcomp[0.10]/across_logs_best_alternative.tex"}}
\newcommand{\AppendixSensitivityConfigurationsInput}{\input{"../../../results/perm/powerlaw_statistics/real/sensitivity/configurations.tex"}}
\newcommand{\AppendixSensitivitySummaryInput}{\input{"../../../results/perm/powerlaw_statistics/real/sensitivity/sensitivity_summary.tex"}}
\newcommand{\AppendixExploreDirPath}{\texttt{results/perm/powerlaw\_statistics/real/xmin-adaptive/explore\_gof\_changes/}}
\newcommand{\AppendixClausetConfigDirPath}{\texttt{results/perm/powerlaw\_statistics/real/xmin-adaptive/clauset\_results/$<$config$>$/}}

\maketitle

\noindent
This online appendix consolidates supporting empirical tables for the ICPM 2027 submission \textit{Assessing Power-Law Compatibility of Process Behavior Frequencies Across Orders of Behavioral Abstraction}.
The supporting data stems from the paper's online repository and is compiled into this PDF document for easier viewing.
All provided file references are relevative for the code repository.

% Depth 3 includes Section~4.3 sensitivity configuration headings.
\setcounter{tocdepth}{3}
\renewcommand{\contentsname}{Contents}
{\small\tableofcontents}

\section{Event Logs and Preprocessing}
\label{sec:preprocessing}

\subsection{Event logs}
\label{sec:event-logs}

The analysis uses the following real-world event logs:

\begin{itemize}\footnotesize\raggedright
\item \textbf{ACCRE (Academic Credentialing)}: Chapela-Campa, D., Benchekroun, I., Baron, O., Dumas, M., Krass, D., \& Senderovich, A. (2024). Evaluation datasets and results of the paper ``A Framework for Measuring the Quality of Business Process Simulation Models'' [Data set]. Zenodo. \url{https://doi.org/10.5281/zenodo.12126071}
\item \textbf{BPIC11 (BPI Challenge 2011, Dutch Academic Hospital)}: van Dongen, B. (2011). Real-life event logs - Hospital log (Version 1) [Data set]. 4TU.ResearchData. \url{https://doi.org/10.4121/uuid:d9769f3d-0ab0-4fb8-803b-0d1120ffcf54}
\item \textbf{BPIC12 (BPI Challenge 2012)}: van Dongen, B. (2012). BPI Challenge 2012 (Version 1) [Data set]. 4TU.ResearchData. \url{https://doi.org/10.4121/uuid:3926db30-f712-4394-aebc-75976070e91f}
\item \textbf{BPIC13\_cp (BPI Challenge 2013, closed problems)}: Steeman, W. (2013). BPI Challenge 2013, closed problems (Version 1) [Data set]. 4TU.ResearchData. \url{https://doi.org/10.4121/uuid:c2c3b154-ab26-4b31-a0e8-8f2350ddac11}
\item \textbf{BPIC13\_i (BPI Challenge 2013, incidents)}: Steeman, W. (2013). BPI Challenge 2013, incidents (Version 1) [Data set]. 4TU.ResearchData. \url{https://doi.org/10.4121/uuid:500573e6-accc-4b0c-9576-aa5468b10cee}
\item \textbf{BPIC14 (BPI Challenge 2014)}: van Dongen, B. (2014). BPI Challenge 2014: Incident details (Version 1) [Data set]. 4TU.ResearchData. \url{https://doi.org/10.4121/uuid:3cfa2260-f5c5-44be-afe1-b70d35288d6d}
\item \textbf{BPIC15\_1 (BPI Challenge 2015, Municipality 1)}: van Dongen, B.~F. (2015). BPI Challenge 2015 (Version 1) [Data set]. 4TU.ResearchData. \url{https://doi.org/10.4121/uuid:31a308ef-c844-48da-948c-305d167a0ec1}
\item \textbf{BPIC15\_2 (BPI Challenge 2015, Municipality 2)}: van Dongen, B.~F. (2015). BPI Challenge 2015 (Version 1) [Data set]. 4TU.ResearchData. \url{https://doi.org/10.4121/uuid:31a308ef-c844-48da-948c-305d167a0ec1}
\item \textbf{BPIC15\_3 (BPI Challenge 2015, Municipality 3)}: van Dongen, B.~F. (2015). BPI Challenge 2015 (Version 1) [Data set]. 4TU.ResearchData. \url{https://doi.org/10.4121/uuid:31a308ef-c844-48da-948c-305d167a0ec1}
\item \textbf{BPIC15\_4 (BPI Challenge 2015, Municipality 4)}: van Dongen, B.~F. (2015). BPI Challenge 2015 (Version 1) [Data set]. 4TU.ResearchData. \url{https://doi.org/10.4121/uuid:31a308ef-c844-48da-948c-305d167a0ec1}
\item \textbf{BPIC15\_5 (BPI Challenge 2015, Municipality 5)}: van Dongen, B.~F. (2015). BPI Challenge 2015 (Version 1) [Data set]. 4TU.ResearchData. \url{https://doi.org/10.4121/uuid:31a308ef-c844-48da-948c-305d167a0ec1}
\item \textbf{BPIC17 (BPI Challenge 2017)}: van Dongen, B. (2017). BPI Challenge 2017 (Version 1) [Data set]. 4TU.ResearchData. \url{https://doi.org/10.4121/uuid:5f3067df-f10b-45da-b98b-86ae4c7a310b}
\item \textbf{BPIC18 (BPI Challenge 2018, EU Agriculture Subsidy)}: van Dongen, B., \& Borchert, F. (2018). BPI Challenge 2018 (Version 1) [Data set]. 4TU.ResearchData. \url{https://doi.org/10.4121/uuid:3301445f-95e8-4ff0-98a4-901f1f204972}
\item \textbf{BPIC19 (BPI Challenge 2019)}: van Dongen, B. (2019). BPI Challenge 2019 (Version 1) [Data set]. 4TU.ResearchData. \url{https://doi.org/10.4121/uuid:d06aff4b-79f0-45e6-8ec8-e19730c248f1}
\item \textbf{BPIC20\_dd (BPI Challenge 2020, Domestic Declarations)}: van Dongen, B. (2020). BPI Challenge 2020: Domestic Declarations (Version 1) [Data set]. 4TU.ResearchData. \url{https://doi.org/10.4121/uuid:3f422315-ed9d-4882-891f-e180b5b4feb5}
\item \textbf{BPIC20\_id (BPI Challenge 2020, International Declarations)}: van Dongen, B. (2020). BPI Challenge 2020: International Declarations (Version 1) [Data set]. 4TU.ResearchData. \url{https://doi.org/10.4121/uuid:2bbf8f6a-fc50-48eb-aa9e-c4ea5ef7e8c5}
\item \textbf{BPIC20\_ptc (BPI Challenge 2020, Prepaid Travel Costs)}: van Dongen, B. (2020). BPI Challenge 2020: Prepaid Travel Costs (Version 1) [Data set]. 4TU.ResearchData. \url{https://doi.org/10.4121/uuid:5d2fe5e1-f91f-4a3b-ad9b-9e4126870165}
\item \textbf{BPIC20\_rfp (BPI Challenge 2020, Request For Payment)}: van Dongen, B. (2020). BPI Challenge 2020: Request For Payment (Version 1) [Data set]. 4TU.ResearchData. \url{https://doi.org/10.4121/uuid:895b26fb-6f25-46eb-9e48-0dca26fcd030}
\item \textbf{BPIC20\_tpd (BPI Challenge 2020, Travel Permit Data)}: van Dongen, B. (2020). BPI Challenge 2020: Travel Permit Data (Version 1) [Data set]. 4TU.ResearchData. \url{https://doi.org/10.4121/uuid:ea03d361-a7cd-4f5e-83d8-5fbdf0362550}
\item \textbf{CALL}: Chapela-Campa, D., Benchekroun, I., Baron, O., Dumas, M., Krass, D., \& Senderovich, A. (2024). Evaluation datasets and results of the paper ``A Framework for Measuring the Quality of Business Process Simulation Models'' [Data set]. Zenodo. \url{https://doi.org/10.5281/zenodo.12126071}
\item \textbf{ITHD (Help Desk)}: Polato, M. (2017). Dataset belonging to the help desk log of an Italian Company (Version 1) [Data set]. 4TU.ResearchData. \url{https://doi.org/10.4121/uuid:0c60edf1-6f83-4e75-9367-4c63b3e9d5bb}
\item \textbf{MOBIS (MobIS Challenge 2019)}: Scheid, M., Rehse, J.-R., Houy, C., \& Fettke, P. (2018). Data Set for MobIS Challenge 2019 [Data set]. German Research Center for Artificial Intelligence (DFKI). \url{https://www.dfki.de/web/forschung/projekte-publikationen/publikation/10258}
\item \textbf{RTFMP (Road Traffic Fine Management Process)}: de Leoni, M., \& Mannhardt, F. (2015). Road Traffic Fine Management Process (Version 1) [Data set]. 4TU.ResearchData. \url{https://doi.org/10.4121/uuid:270fd440-1057-4fb9-89a9-b699b47990f5}
\item \textbf{SEPSIS (Sepsis Cases)}: Mannhardt, F. (2016). Sepsis Cases - Event Log (Version 1) [Data set]. 4TU.ResearchData. \url{https://doi.org/10.4121/uuid:915d2bfb-7e84-49ad-a286-dc35f063a460}
\end{itemize}

\subsection{Preprocessing}
\label{sec:preprocessing-details}

Table~\ref{tab:preprocessing} documents the preprocessing choices for the 24 empirical event logs.
The preprocessing is performed from the notebook \path{analyses/01_preprocessing.ipynb} and documented in the repository README.
At analysis time, all logs are loaded via \texttt{utils/io/event\_logs.py}
with ascending \texttt{time:timestamp} sorting; equal timestamps retain XES
parse order.
Activity labels follow \texttt{utils/io/activity\_labels.py}: if an XES
classifier named \texttt{Activity classifier} is present, its keys are joined
with \texttt{+}. Otherwise the activity name (\texttt{concept:name}) is used.

\begin{table}[H]
\centering
\caption{Preprocessing choices for analyzed real event logs.}
\label{tab:preprocessing}
\resizebox{\textwidth}{!}{%
\begin{tabular}{@{}lllll@{}}
\toprule
Log & Case id & Activity attribute & Timestamp & Filtering \\
\midrule
ACCRE & \texttt{case\_id} & \texttt{activity} & \texttt{start\_time} & Train split only \\
BPIC11 & trace \texttt{concept:name} & \texttt{concept:name} & \texttt{time:timestamp} & None \\
BPIC12 & trace \texttt{concept:name} & name + lifecycle (classifier) & \texttt{time:timestamp} & None \\
BPIC13\_cp & trace \texttt{concept:name} & name + lifecycle (classifier) & \texttt{time:timestamp} & None \\
BPIC13\_i & trace \texttt{concept:name} & name + lifecycle (classifier) & \texttt{time:timestamp} & None \\
BPIC14 & \texttt{Incident ID} & \texttt{IncidentActivity\_Type} & \texttt{DateStamp} & None \\
BPIC15\_1 & trace \texttt{concept:name} & name + lifecycle (classifier) & \texttt{time:timestamp} & None \\
BPIC15\_2 & trace \texttt{concept:name} & name + lifecycle (classifier) & \texttt{time:timestamp} & None \\
BPIC15\_3 & trace \texttt{concept:name} & name + lifecycle (classifier) & \texttt{time:timestamp} & None \\
BPIC15\_4 & trace \texttt{concept:name} & name + lifecycle (classifier) & \texttt{time:timestamp} & None \\
BPIC15\_5 & trace \texttt{concept:name} & name + lifecycle (classifier) & \texttt{time:timestamp} & None \\
BPIC17 & trace \texttt{concept:name} & \texttt{concept:name} & \texttt{time:timestamp} & None \\
BPIC18 & trace \texttt{concept:name} & \texttt{concept:name} & \texttt{time:timestamp} & None \\
BPIC19 & trace \texttt{concept:name} & \texttt{concept:name} & \texttt{time:timestamp} & None \\
BPIC20\_dd & trace \texttt{concept:name} & \texttt{concept:name} & \texttt{time:timestamp} & None \\
BPIC20\_id & trace \texttt{concept:name} & \texttt{concept:name} & \texttt{time:timestamp} & None \\
BPIC20\_ptc & trace \texttt{concept:name} & \texttt{concept:name} & \texttt{time:timestamp} & None \\
BPIC20\_rfp & trace \texttt{concept:name} & \texttt{concept:name} & \texttt{time:timestamp} & None \\
BPIC20\_tpd & trace \texttt{concept:name} & \texttt{concept:name} & \texttt{time:timestamp} & None \\
CALL & \texttt{case\_id} & \texttt{activity} & \texttt{end\_time} & Train split only \\
ITHD & \texttt{Case ID} & \texttt{Activity} & \texttt{Complete Timestamp} & None \\
MOBIS & \texttt{case} & \texttt{activity} & \texttt{end} & None \\
RTFMP & trace \texttt{concept:name} & \texttt{concept:name} & \texttt{time:timestamp} & None \\
SEPSIS & trace \texttt{concept:name} & \texttt{concept:name} & \texttt{time:timestamp} & None \\
\bottomrule
\end{tabular}%
}
\end{table}
\FloatBarrier

\section{Alternative Model Comparisons}
\label{sec:alternatives}

For each log, Table~\ref{tab:variant-model-comparisons} lists every comparison of the fitted power law with an alternative distribution. Then, Table~\ref{tab:variant-comparison-summary} counts these outcomes for logs that are not in category A or B. Finally, Table~\ref{tab:best-alternative} reports, for each log and abstraction level, the alternative distribution with the smallest $\mathcal{R}$ relative to the fitted power law.
Accordingly, PL indicates that the power law achieves a higher likelihood than the examined alternatives.
Bold entries indicate statistically significant pairwise comparisons at $p<0.10$.

\AppendixVariantComparisonInput

\AppendixVariantComparisonSummaryInput


\begin{landscape}
\begin{table}[p]
\centering
\caption{Best-fitting distribution relative to the power law across
abstraction orders ($k=1,\ldots,20$) and trace variants (TV).
Bold entries indicate a significant likelihood-ratio comparison
($p<0.10$).}
\label{tab:best-alternative}
\begingroup
\setlength{\tabcolsep}{2pt}%
\tiny
\AppendixBestAlternativeInput
\endgroup
\end{table}
\end{landscape}

\noindent


\section{Reconciliation of Type Counts}
\label{sec:transitions}

The following tables summarise type and occurrence transitions when moving
across abstraction orders for each analysed event log.

% Auto-generated by generate.py — do not edit by hand.
% Deterministic order: ALL_REAL_LOG_DATASETS filtered to available TeX.

\noindent\textit{Log-to-table index}
(click a log name to jump to its reconciliation table):
\begin{center}\footnotesize
\begin{tabular}{@{}*{4}{l}@{}}
\hyperref[tab:transition-ACCRE]{\texttt{ACCRE}}~(p.~\pageref{tab:transition-ACCRE}) & \hyperref[tab:transition-BPIC11]{\texttt{BPIC11}}~(p.~\pageref{tab:transition-BPIC11}) & \hyperref[tab:transition-BPIC12]{\texttt{BPIC12}}~(p.~\pageref{tab:transition-BPIC12}) & \hyperref[tab:transition-BPIC13_cp]{\texttt{BPIC13\_cp}}~(p.~\pageref{tab:transition-BPIC13_cp}) \\
\hyperref[tab:transition-BPIC13_i]{\texttt{BPIC13\_i}}~(p.~\pageref{tab:transition-BPIC13_i}) & \hyperref[tab:transition-BPIC14]{\texttt{BPIC14}}~(p.~\pageref{tab:transition-BPIC14}) & \hyperref[tab:transition-BPIC15_1]{\texttt{BPIC15\_1}}~(p.~\pageref{tab:transition-BPIC15_1}) & \hyperref[tab:transition-BPIC15_2]{\texttt{BPIC15\_2}}~(p.~\pageref{tab:transition-BPIC15_2}) \\
\hyperref[tab:transition-BPIC15_3]{\texttt{BPIC15\_3}}~(p.~\pageref{tab:transition-BPIC15_3}) & \hyperref[tab:transition-BPIC15_4]{\texttt{BPIC15\_4}}~(p.~\pageref{tab:transition-BPIC15_4}) & \hyperref[tab:transition-BPIC15_5]{\texttt{BPIC15\_5}}~(p.~\pageref{tab:transition-BPIC15_5}) & \hyperref[tab:transition-BPIC17]{\texttt{BPIC17}}~(p.~\pageref{tab:transition-BPIC17}) \\
\hyperref[tab:transition-BPIC18]{\texttt{BPIC18}}~(p.~\pageref{tab:transition-BPIC18}) & \hyperref[tab:transition-BPIC19]{\texttt{BPIC19}}~(p.~\pageref{tab:transition-BPIC19}) & \hyperref[tab:transition-BPIC20_dd]{\texttt{BPIC20\_dd}}~(p.~\pageref{tab:transition-BPIC20_dd}) & \hyperref[tab:transition-BPIC20_id]{\texttt{BPIC20\_id}}~(p.~\pageref{tab:transition-BPIC20_id}) \\
\hyperref[tab:transition-BPIC20_ptc]{\texttt{BPIC20\_ptc}}~(p.~\pageref{tab:transition-BPIC20_ptc}) & \hyperref[tab:transition-BPIC20_rfp]{\texttt{BPIC20\_rfp}}~(p.~\pageref{tab:transition-BPIC20_rfp}) & \hyperref[tab:transition-BPIC20_tpd]{\texttt{BPIC20\_tpd}}~(p.~\pageref{tab:transition-BPIC20_tpd}) & \hyperref[tab:transition-CALL]{\texttt{CALL}}~(p.~\pageref{tab:transition-CALL}) \\
\hyperref[tab:transition-ITHD]{\texttt{ITHD}}~(p.~\pageref{tab:transition-ITHD}) & \hyperref[tab:transition-MOBIS]{\texttt{MOBIS}}~(p.~\pageref{tab:transition-MOBIS}) & \hyperref[tab:transition-RTFMP]{\texttt{RTFMP}}~(p.~\pageref{tab:transition-RTFMP}) & \hyperref[tab:transition-SEPSIS]{\texttt{SEPSIS}}~(p.~\pageref{tab:transition-SEPSIS}) \\
\end{tabular}
\end{center}

\begin{landscape}
\begin{table}[p]
\centering
\caption{Type-count reconciliation for \texttt{ACCRE}.}
\label{tab:transition-ACCRE}
\input{"../../../results/pl_stats/real/xmin-adaptive/explore_gof_changes/ACCRE/transition_overview.tex"}
\end{table}
\end{landscape}

\begin{landscape}
\begin{table}[p]
\centering
\caption{Type-count reconciliation for \texttt{BPIC11}.}
\label{tab:transition-BPIC11}
\input{"../../../results/pl_stats/real/xmin-adaptive/explore_gof_changes/BPIC11/transition_overview.tex"}
\end{table}
\end{landscape}

\begin{landscape}
\begin{table}[p]
\centering
\caption{Type-count reconciliation for \texttt{BPIC12}.}
\label{tab:transition-BPIC12}
\input{"../../../results/pl_stats/real/xmin-adaptive/explore_gof_changes/BPIC12/transition_overview.tex"}
\end{table}
\end{landscape}

\begin{landscape}
\begin{table}[p]
\centering
\caption{Type-count reconciliation for \texttt{BPIC13\_cp}.}
\label{tab:transition-BPIC13_cp}
\input{"../../../results/pl_stats/real/xmin-adaptive/explore_gof_changes/BPIC13_cp/transition_overview.tex"}
\end{table}
\end{landscape}

\begin{landscape}
\begin{table}[p]
\centering
\caption{Type-count reconciliation for \texttt{BPIC13\_i}.}
\label{tab:transition-BPIC13_i}
\input{"../../../results/pl_stats/real/xmin-adaptive/explore_gof_changes/BPIC13_i/transition_overview.tex"}
\end{table}
\end{landscape}

\begin{landscape}
\begin{table}[p]
\centering
\caption{Type-count reconciliation for \texttt{BPIC14}.}
\label{tab:transition-BPIC14}
\input{"../../../results/pl_stats/real/xmin-adaptive/explore_gof_changes/BPIC14/transition_overview.tex"}
\end{table}
\end{landscape}

\begin{landscape}
\begin{table}[p]
\centering
\caption{Type-count reconciliation for \texttt{BPIC15\_1}.}
\label{tab:transition-BPIC15_1}
\input{"../../../results/pl_stats/real/xmin-adaptive/explore_gof_changes/BPIC15_1/transition_overview.tex"}
\end{table}
\end{landscape}

\begin{landscape}
\begin{table}[p]
\centering
\caption{Type-count reconciliation for \texttt{BPIC15\_2}.}
\label{tab:transition-BPIC15_2}
\input{"../../../results/pl_stats/real/xmin-adaptive/explore_gof_changes/BPIC15_2/transition_overview.tex"}
\end{table}
\end{landscape}

\begin{landscape}
\begin{table}[p]
\centering
\caption{Type-count reconciliation for \texttt{BPIC15\_3}.}
\label{tab:transition-BPIC15_3}
\input{"../../../results/pl_stats/real/xmin-adaptive/explore_gof_changes/BPIC15_3/transition_overview.tex"}
\end{table}
\end{landscape}

\begin{landscape}
\begin{table}[p]
\centering
\caption{Type-count reconciliation for \texttt{BPIC15\_4}.}
\label{tab:transition-BPIC15_4}
\input{"../../../results/pl_stats/real/xmin-adaptive/explore_gof_changes/BPIC15_4/transition_overview.tex"}
\end{table}
\end{landscape}

\begin{landscape}
\begin{table}[p]
\centering
\caption{Type-count reconciliation for \texttt{BPIC15\_5}.}
\label{tab:transition-BPIC15_5}
\input{"../../../results/pl_stats/real/xmin-adaptive/explore_gof_changes/BPIC15_5/transition_overview.tex"}
\end{table}
\end{landscape}

\begin{landscape}
\begin{table}[p]
\centering
\caption{Type-count reconciliation for \texttt{BPIC17}.}
\label{tab:transition-BPIC17}
\input{"../../../results/pl_stats/real/xmin-adaptive/explore_gof_changes/BPIC17/transition_overview.tex"}
\end{table}
\end{landscape}

\begin{landscape}
\begin{table}[p]
\centering
\caption{Type-count reconciliation for \texttt{BPIC18}.}
\label{tab:transition-BPIC18}
\input{"../../../results/pl_stats/real/xmin-adaptive/explore_gof_changes/BPIC18/transition_overview.tex"}
\end{table}
\end{landscape}

\begin{landscape}
\begin{table}[p]
\centering
\caption{Type-count reconciliation for \texttt{BPIC19}.}
\label{tab:transition-BPIC19}
\input{"../../../results/pl_stats/real/xmin-adaptive/explore_gof_changes/BPIC19/transition_overview.tex"}
\end{table}
\end{landscape}

\begin{landscape}
\begin{table}[p]
\centering
\caption{Type-count reconciliation for \texttt{BPIC20\_dd}.}
\label{tab:transition-BPIC20_dd}
\input{"../../../results/pl_stats/real/xmin-adaptive/explore_gof_changes/BPIC20_dd/transition_overview.tex"}
\end{table}
\end{landscape}

\begin{landscape}
\begin{table}[p]
\centering
\caption{Type-count reconciliation for \texttt{BPIC20\_id}.}
\label{tab:transition-BPIC20_id}
\input{"../../../results/pl_stats/real/xmin-adaptive/explore_gof_changes/BPIC20_id/transition_overview.tex"}
\end{table}
\end{landscape}

\begin{landscape}
\begin{table}[p]
\centering
\caption{Type-count reconciliation for \texttt{BPIC20\_ptc}.}
\label{tab:transition-BPIC20_ptc}
\input{"../../../results/pl_stats/real/xmin-adaptive/explore_gof_changes/BPIC20_ptc/transition_overview.tex"}
\end{table}
\end{landscape}

\begin{landscape}
\begin{table}[p]
\centering
\caption{Type-count reconciliation for \texttt{BPIC20\_rfp}.}
\label{tab:transition-BPIC20_rfp}
\input{"../../../results/pl_stats/real/xmin-adaptive/explore_gof_changes/BPIC20_rfp/transition_overview.tex"}
\end{table}
\end{landscape}

\begin{landscape}
\begin{table}[p]
\centering
\caption{Type-count reconciliation for \texttt{BPIC20\_tpd}.}
\label{tab:transition-BPIC20_tpd}
\input{"../../../results/pl_stats/real/xmin-adaptive/explore_gof_changes/BPIC20_tpd/transition_overview.tex"}
\end{table}
\end{landscape}

\begin{landscape}
\begin{table}[p]
\centering
\caption{Type-count reconciliation for \texttt{CALL}.}
\label{tab:transition-CALL}
\input{"../../../results/pl_stats/real/xmin-adaptive/explore_gof_changes/CALL/transition_overview.tex"}
\end{table}
\end{landscape}

\begin{landscape}
\begin{table}[p]
\centering
\caption{Type-count reconciliation for \texttt{ITHD}.}
\label{tab:transition-ITHD}
\input{"../../../results/pl_stats/real/xmin-adaptive/explore_gof_changes/ITHD/transition_overview.tex"}
\end{table}
\end{landscape}

\begin{landscape}
\begin{table}[p]
\centering
\caption{Type-count reconciliation for \texttt{MOBIS}.}
\label{tab:transition-MOBIS}
\input{"../../../results/pl_stats/real/xmin-adaptive/explore_gof_changes/MOBIS/transition_overview.tex"}
\end{table}
\end{landscape}

\begin{landscape}
\begin{table}[p]
\centering
\caption{Type-count reconciliation for \texttt{RTFMP}.}
\label{tab:transition-RTFMP}
\input{"../../../results/pl_stats/real/xmin-adaptive/explore_gof_changes/RTFMP/transition_overview.tex"}
\end{table}
\end{landscape}

\begin{landscape}
\begin{table}[p]
\centering
\caption{Type-count reconciliation for \texttt{SEPSIS}.}
\label{tab:transition-SEPSIS}
\input{"../../../results/pl_stats/real/xmin-adaptive/explore_gof_changes/SEPSIS/transition_overview.tex"}
\end{table}
\end{landscape}



\section{Sensitivity Analysis}
\label{sec:sensitivity}

Threshold configurations are varied one-at-a-time. Fits, GOF bootstraps, and likelihood-ratio tests are not recomputed as these do not depend on the configuration.

\subsection{Configuration overview}
\label{sec:sensitivity-config}

\begin{table}[H]
\centering
\caption{One-at-a-time sensitivity configurations used in the appendix.}
\label{tab:sensitivity-configurations}
\AppendixSensitivityConfigurationsInput
\end{table}

\subsection{Summary across configurations}
\label{sec:sensitivity-summary}

\begin{table}[htbp]
\centering
\caption{Summary of trace-variant outcomes and paired GOF diagnostics across
sensitivity configurations.}
\label{tab:sensitivity-summary}
\begingroup
\setlength{\tabcolsep}{2.5pt}%
\scriptsize
\AppendixSensitivitySummaryInput
\endgroup
\end{table}

\subsection{Configuration-specific results}
\label{sec:sensitivity-detail}

For each sensitivity configuration, we include the trace-variant power-law
table, the across-logs category matrix, and the paired GOF comparison figure
written under
\AppendixClausetConfigDirPath{}.

% Auto-generated by generate.py — do not edit by hand.
% Deterministic order follows sensitivity.CONFIGURATIONS.

\subsubsection{Baseline}
\label{sec:sens-baseline}
\FloatBarrier

\begin{table}[H]
\centering
\caption{Power-law analysis results per log at variant level for configuration \texttt{baseline}.}
\label{tab:sens-variant-baseline}
\setlength{\tabcolsep}{2pt}%
\input{"../../../results/pl_stats/real/xmin-adaptive/clauset_results/min[50]-gofp[0.10]-pcomp[0.10]/variant_power_law.tex"}
\end{table}

\begin{table}[H]
\centering
\caption{Power-law analysis results at increasing abstraction levels for configuration \texttt{baseline}.}
\label{tab:sens-categories-baseline}
\setlength{\tabcolsep}{2pt}%
\input{"../../../results/pl_stats/real/xmin-adaptive/clauset_results/min[50]-gofp[0.10]-pcomp[0.10]/across_logs_categories.tex"}
\end{table}

\begin{figure}[H]
\centering
\includegraphics[width=\textwidth]{"../../../results/pl_stats/real/xmin-adaptive/clauset_results/min[50]-gofp[0.10]-pcomp[0.10]/paired_gof_comparison.pdf"}
\caption{Paired differences in goodness-of-fit non-rejection shares (trace variants minus finite order) with pointwise 95\% confidence intervals for configuration \texttt{baseline}.}
\label{fig:sens-paired-baseline}
\end{figure}
\FloatBarrier

\subsubsection{Minimum fitted types = 0}
\label{sec:sens-min_types_0}
\FloatBarrier

\begin{table}[H]
\centering
\caption{Power-law analysis results per log at variant level for configuration \texttt{min\_types\_0}.}
\label{tab:sens-variant-min_types_0}
\setlength{\tabcolsep}{2pt}%
\input{"../../../results/pl_stats/real/xmin-adaptive/clauset_results/min[00]-gofp[0.10]-pcomp[0.10]/variant_power_law.tex"}
\end{table}

\begin{table}[H]
\centering
\caption{Power-law analysis results at increasing abstraction levels for configuration \texttt{min\_types\_0}.}
\label{tab:sens-categories-min_types_0}
\setlength{\tabcolsep}{2pt}%
\input{"../../../results/pl_stats/real/xmin-adaptive/clauset_results/min[00]-gofp[0.10]-pcomp[0.10]/across_logs_categories.tex"}
\end{table}

\begin{figure}[H]
\centering
\includegraphics[width=\textwidth]{"../../../results/pl_stats/real/xmin-adaptive/clauset_results/min[00]-gofp[0.10]-pcomp[0.10]/paired_gof_comparison.pdf"}
\caption{Paired differences in goodness-of-fit non-rejection shares (trace variants minus finite order) with pointwise 95\% confidence intervals for configuration \texttt{min\_types\_0}.}
\label{fig:sens-paired-min_types_0}
\end{figure}
\FloatBarrier

\subsubsection{Minimum fitted types = 15}
\label{sec:sens-min_types_15}
\FloatBarrier

\begin{table}[H]
\centering
\caption{Power-law analysis results per log at variant level for configuration \texttt{min\_types\_15}.}
\label{tab:sens-variant-min_types_15}
\setlength{\tabcolsep}{2pt}%
\input{"../../../results/pl_stats/real/xmin-adaptive/clauset_results/min[15]-gofp[0.10]-pcomp[0.10]/variant_power_law.tex"}
\end{table}

\begin{table}[H]
\centering
\caption{Power-law analysis results at increasing abstraction levels for configuration \texttt{min\_types\_15}.}
\label{tab:sens-categories-min_types_15}
\setlength{\tabcolsep}{2pt}%
\input{"../../../results/pl_stats/real/xmin-adaptive/clauset_results/min[15]-gofp[0.10]-pcomp[0.10]/across_logs_categories.tex"}
\end{table}

\begin{figure}[H]
\centering
\includegraphics[width=\textwidth]{"../../../results/pl_stats/real/xmin-adaptive/clauset_results/min[15]-gofp[0.10]-pcomp[0.10]/paired_gof_comparison.pdf"}
\caption{Paired differences in goodness-of-fit non-rejection shares (trace variants minus finite order) with pointwise 95\% confidence intervals for configuration \texttt{min\_types\_15}.}
\label{fig:sens-paired-min_types_15}
\end{figure}
\FloatBarrier

\subsubsection{Minimum fitted types = 30}
\label{sec:sens-min_types_30}
\FloatBarrier

\begin{table}[H]
\centering
\caption{Power-law analysis results per log at variant level for configuration \texttt{min\_types\_30}.}
\label{tab:sens-variant-min_types_30}
\setlength{\tabcolsep}{2pt}%
\input{"../../../results/pl_stats/real/xmin-adaptive/clauset_results/min[30]-gofp[0.10]-pcomp[0.10]/variant_power_law.tex"}
\end{table}

\begin{table}[H]
\centering
\caption{Power-law analysis results at increasing abstraction levels for configuration \texttt{min\_types\_30}.}
\label{tab:sens-categories-min_types_30}
\setlength{\tabcolsep}{2pt}%
\input{"../../../results/pl_stats/real/xmin-adaptive/clauset_results/min[30]-gofp[0.10]-pcomp[0.10]/across_logs_categories.tex"}
\end{table}

\begin{figure}[H]
\centering
\includegraphics[width=\textwidth]{"../../../results/pl_stats/real/xmin-adaptive/clauset_results/min[30]-gofp[0.10]-pcomp[0.10]/paired_gof_comparison.pdf"}
\caption{Paired differences in goodness-of-fit non-rejection shares (trace variants minus finite order) with pointwise 95\% confidence intervals for configuration \texttt{min\_types\_30}.}
\label{fig:sens-paired-min_types_30}
\end{figure}
\FloatBarrier

\subsubsection{GOF threshold = 0.05}
\label{sec:sens-gof_005}
\FloatBarrier

\begin{table}[H]
\centering
\caption{Power-law analysis results per log at variant level for configuration \texttt{gof\_005}.}
\label{tab:sens-variant-gof_005}
\setlength{\tabcolsep}{2pt}%
\input{"../../../results/pl_stats/real/xmin-adaptive/clauset_results/min[50]-gofp[0.05]-pcomp[0.10]/variant_power_law.tex"}
\end{table}

\begin{table}[H]
\centering
\caption{Power-law analysis results at increasing abstraction levels for configuration \texttt{gof\_005}.}
\label{tab:sens-categories-gof_005}
\setlength{\tabcolsep}{2pt}%
\input{"../../../results/pl_stats/real/xmin-adaptive/clauset_results/min[50]-gofp[0.05]-pcomp[0.10]/across_logs_categories.tex"}
\end{table}

\begin{figure}[H]
\centering
\includegraphics[width=\textwidth]{"../../../results/pl_stats/real/xmin-adaptive/clauset_results/min[50]-gofp[0.05]-pcomp[0.10]/paired_gof_comparison.pdf"}
\caption{Paired differences in goodness-of-fit non-rejection shares (trace variants minus finite order) with pointwise 95\% confidence intervals for configuration \texttt{gof\_005}.}
\label{fig:sens-paired-gof_005}
\end{figure}
\FloatBarrier

\subsubsection{GOF threshold = 0.20}
\label{sec:sens-gof_020}
\FloatBarrier

\begin{table}[H]
\centering
\caption{Power-law analysis results per log at variant level for configuration \texttt{gof\_020}.}
\label{tab:sens-variant-gof_020}
\setlength{\tabcolsep}{2pt}%
\input{"../../../results/pl_stats/real/xmin-adaptive/clauset_results/min[50]-gofp[0.20]-pcomp[0.10]/variant_power_law.tex"}
\end{table}

\begin{table}[H]
\centering
\caption{Power-law analysis results at increasing abstraction levels for configuration \texttt{gof\_020}.}
\label{tab:sens-categories-gof_020}
\setlength{\tabcolsep}{2pt}%
\input{"../../../results/pl_stats/real/xmin-adaptive/clauset_results/min[50]-gofp[0.20]-pcomp[0.10]/across_logs_categories.tex"}
\end{table}

\begin{figure}[H]
\centering
\includegraphics[width=\textwidth]{"../../../results/pl_stats/real/xmin-adaptive/clauset_results/min[50]-gofp[0.20]-pcomp[0.10]/paired_gof_comparison.pdf"}
\caption{Paired differences in goodness-of-fit non-rejection shares (trace variants minus finite order) with pointwise 95\% confidence intervals for configuration \texttt{gof\_020}.}
\label{fig:sens-paired-gof_020}
\end{figure}
\FloatBarrier

\subsubsection{Comparison threshold = 0.05}
\label{sec:sens-comp_005}
\FloatBarrier

\begin{table}[H]
\centering
\caption{Power-law analysis results per log at variant level for configuration \texttt{comp\_005}.}
\label{tab:sens-variant-comp_005}
\setlength{\tabcolsep}{2pt}%
\input{"../../../results/pl_stats/real/xmin-adaptive/clauset_results/min[50]-gofp[0.10]-pcomp[0.05]/variant_power_law.tex"}
\end{table}

\begin{table}[H]
\centering
\caption{Power-law analysis results at increasing abstraction levels for configuration \texttt{comp\_005}.}
\label{tab:sens-categories-comp_005}
\setlength{\tabcolsep}{2pt}%
\input{"../../../results/pl_stats/real/xmin-adaptive/clauset_results/min[50]-gofp[0.10]-pcomp[0.05]/across_logs_categories.tex"}
\end{table}

\begin{figure}[H]
\centering
\includegraphics[width=\textwidth]{"../../../results/pl_stats/real/xmin-adaptive/clauset_results/min[50]-gofp[0.10]-pcomp[0.05]/paired_gof_comparison.pdf"}
\caption{Paired differences in goodness-of-fit non-rejection shares (trace variants minus finite order) with pointwise 95\% confidence intervals for configuration \texttt{comp\_005}.}
\label{fig:sens-paired-comp_005}
\end{figure}
\FloatBarrier

\subsubsection{Comparison threshold = 0.20}
\label{sec:sens-comp_020}
\FloatBarrier

\begin{table}[H]
\centering
\caption{Power-law analysis results per log at variant level for configuration \texttt{comp\_020}.}
\label{tab:sens-variant-comp_020}
\setlength{\tabcolsep}{2pt}%
\input{"../../../results/pl_stats/real/xmin-adaptive/clauset_results/min[50]-gofp[0.10]-pcomp[0.20]/variant_power_law.tex"}
\end{table}

\begin{table}[H]
\centering
\caption{Power-law analysis results at increasing abstraction levels for configuration \texttt{comp\_020}.}
\label{tab:sens-categories-comp_020}
\setlength{\tabcolsep}{2pt}%
\input{"../../../results/pl_stats/real/xmin-adaptive/clauset_results/min[50]-gofp[0.10]-pcomp[0.20]/across_logs_categories.tex"}
\end{table}

\begin{figure}[H]
\centering
\includegraphics[width=\textwidth]{"../../../results/pl_stats/real/xmin-adaptive/clauset_results/min[50]-gofp[0.10]-pcomp[0.20]/paired_gof_comparison.pdf"}
\caption{Paired differences in goodness-of-fit non-rejection shares (trace variants minus finite order) with pointwise 95\% confidence intervals for configuration \texttt{comp\_020}.}
\label{fig:sens-paired-comp_020}
\end{figure}
\FloatBarrier



\section{Simulation Use Case Example}
\label{sec:simulation-example}

This section reports the trace-variant assessment of the synthetic CVS pharmacy event log from Camargo et al.\ (2022). The log is available at \url{https://zenodo.org/records/5734443}. Table~\ref{tab:synth-cvs-min50} reports the results using the baseline minimum of 50 fitted types. Table~\ref{tab:synth-cvs-min30} reports the results after lowering this minimum to 30 fitted types.

\begin{quote}
Camargo, M., Dumas, M., \& Gonz\'alez Rojas, O. (2022).
Learning Accurate Business Process Simulation Models from Event Logs via Automated Process Discovery and Deep Learning.
In \textit{CAiSE 2022} (pp.~55--71). Springer.
\end{quote}

\begin{table}[H]
\centering
\caption{Power-law analysis results per log at variant level for the synthetic
CVS pharmacy log, with a minimum of 50 fitted types.}
\label{tab:synth-cvs-min50}
\setlength{\tabcolsep}{2pt}%
\input{"../../../results/perm/powerlaw_statistics/synth_example/xmin-adaptive/clauset_results/min[50]-gofp[0.10]-pcomp[0.10]/variant_power_law.tex"}
\end{table}

\begin{table}[H]
\centering
\caption{Power-law analysis results per log at variant level for the synthetic
CVS pharmacy log, with a minimum of 30 fitted types.}
\label{tab:synth-cvs-min30}
\setlength{\tabcolsep}{2pt}%
\input{"../../../results/perm/powerlaw_statistics/synth_example/xmin-adaptive/clauset_results/min[30]-gofp[0.10]-pcomp[0.10]/variant_power_law.tex"}
\end{table}
\FloatBarrier

\end{document}
